\chapter{Podstawy algebry abstrakcyjnej}
\label{ch:algebra-abstrakcyjna}
\index{algebra abstrakcyjna@Algebra abstrakcyjna}

Zbierzemy pojęcia grup, pierścieni i modułów potrzebne
do algebry homologicznej. Konstrukcje ilorazowe pozwalają
narzucać relacje, a iloczyn tensorowy pozwala zastąpić
odwzorowanie dwuliniowe liniowym. Wcześniejszy
Rozdział~\ref{ch:algebra-tensorowa} dotyczył przestrzeni
nad ciałem; tutaj pracujemy nad pierścieniem.
W Rozdziale~\ref{ch:grass-spin} wystąpił też moduł
algebry Clifforda; wrócimy do tego przykładu.

\section{Grupy, słowa i prezentacje}
\label{sec:grupy-prezentacje-5}
\index{grupa wolna@Grupa wolna}
\index{prezentacja grupy@Prezentacja grupy}

Zanim przejdziemy do modułów, zbierzmy język
potrzebny przy grupie podstawowej i
Twierdzeniu~\ref{thm:van-kampen}. Grupa opisuje
odwracalne składanie działań. W grupie klas
pętli opartych w jednym punkcie działaniem
jest składanie pętli, a odwrotność reprezentuje
pętla przebiegnięta w przeciwnym kierunku.

\begin{definicja}[Grupa i homomorfizm]
\index{grupa i homomorfizm@Grupa i homomorfizm}
\label{def:grupa-hom-5}
\emph{Grupa} $(G,\cdot,1)$ to zbiór z łącznym
mnożeniem, elementem neutralnym $1$ i
odwrotnością $g^{-1}$ każdego $g\in G$.
Nie zakładamy przemienności. Odwzorowanie
$\varphi:G\to H$ jest \emph{homomorfizmem},
jeśli $\varphi(gh)=\varphi(g)\varphi(h)$.
Wtedy $\varphi(1)=1$ i
$\varphi(g^{-1})=\varphi(g)^{-1}$.
\end{definicja}

Przykładem jest $(\Z,+,0)$;
w zapisie addytywnym odwrotnością $m$
jest $-m$. Grupa permutacji trzech
elementów dostarcza przykładu
nieprzemiennego: przestawienia
$(12)(23)$ i $(23)(12)$ są różne.

\begin{definicja}[Podgrupa normalna i iloraz]
\index{podgrupa normalna i iloraz@Podgrupa normalna i iloraz}
\label{def:normal-quotient-5}
Podgrupa $N\leq G$ jest \emph{normalna},
gdy $gNg^{-1}=N$ dla każdego $g\in G$.
Wówczas zbiór warstw $G/N$ staje się
grupą przez $(gN)(hN)=ghN$.
Jądro homomorfizmu
$\ker\varphi=\{g:\varphi(g)=1\}$
jest normalne. \emph{Domknięcie normalne}
zbioru $R\subset G$, oznaczane
$\langle\!\langle R\rangle\!\rangle$,
jest najmniejszą podgrupą normalną
zawierającą $R$.
\end{definicja}

\begin{proof}[Dlaczego iloraz jest dobrze określony?]
Jeśli $g'=gn$ i $h'=hm$ dla $n,m\in N$,
to $g'h'=gnhm=gh(h^{-1}nh)m$.
Normalność daje $h^{-1}nh\in N$,
zatem $g'h'N=ghN$. Mnożenie
nie zależy więc od wybranych
reprezentantów. Dla jądra:
$\varphi(gng^{-1})=
\varphi(g)1\varphi(g)^{-1}=1$.
Elementy domknięcia normalnego
są skończonymi iloczynami
sprzężeń $grg^{-1}$ i ich
odwrotności, $r\in R$:
zbiór tych iloczynów jest
podgrupą normalną zawierającą
$R$ i mieści się w każdej takiej
podgrupie.
\end{proof}

\begin{przyklad}[Reszty modulo $m$]
\label{prz:Zmodm-grupa-5}
Niech $m\geq2$.
W grupie addytywnej $\Z$ każda
podgrupa jest normalna.
Iloraz $\Z/m\Z$ ma elementy
$[0],\ldots,[m-1]$ i dodawanie
$[a]+[b]=[a+b]$. Homomorfizm
$\Z\to\Z/m\Z$ ma jądro $m\Z$.
\end{przyklad}

\begin{definicja}[Grupa wolna i iloczyn wolny]
\index{grupa wolna i iloczyn wolny@Grupa wolna i iloczyn wolny}
\index{slowo w grupie wolnej@słowo w grupie wolnej}
\label{def:free-groups-5}
Niech $A$ będzie zbiorem \emph{symboli},
zwanych generatorami. Dla każdego $a\in A$
tworzymy nowy symbol $a^{-1}$;
zbiór tych nowych symboli oznaczamy $A^{-1}$
i przyjmujemy $(a^{-1})^{-1}=a$.
\emph{Słowo} w alfabecie
$A^{\pm1}=A\sqcup A^{-1}$ jest skończonym
ciągiem liter $x_1\cdots x_r$, $r\geq0$.
Długość słowa to $r$; jedyne słowo długości
zero oznaczamy $\varnothing$ i nazywamy pustym.
Słowa można sklejać, czyli konkatenować:
$(x_1\cdots x_r)(y_1\cdots y_s)
=x_1\cdots x_ry_1\cdots y_s$.

\emph{Redukcja elementarna} usuwa sąsiednie
litery $xx^{-1}$ lub $x^{-1}x$.
Słowo jest \emph{zredukowane}, gdy nie zawiera
takiej pary. Dwa słowa uznajemy za równoważne,
jeżeli można przejść od jednego do drugiego
skończonym ciągiem wstawień lub usunięć
takich par. \emph{Grupa wolna} $F(A)$
to zbiór tych klas, z mnożeniem zadanym
przez sklejanie słów.

Dla grup $G,H$ \emph{iloczyn wolny} $G*H$
powstaje analogicznie ze słów o literach
z $G\setminus\{1_G\}$ i $H\setminus\{1_H\}$.
Redukcja mnoży sąsiednie litery pochodzące
z tego samego czynnika i usuwa otrzymaną
jedynkę. Słowo zredukowane ma litery
z $G$ i $H$ naprzemiennie.
\end{definicja}

Zanim utożsamimy klasy ze słowami zredukowanymi,
trzeba wiedzieć, że wynik redukcji nie zależy
od kolejności kasowania par. Sam zbiór
\emph{wszystkich} słów z konkatenacją
jest tylko monoidem: słowa $aa^{-1}$
i $\varnothing$ nie są w nim dosłownie równe.

\begin{stwierdzenie}[Jednoznaczna postać zredukowana]
\index{redukcja slowa@redukcja słowa}
\label{prop:free-reduction-5}
Każda klasa w $F(A)$ zawiera dokładnie jedno
słowo zredukowane. Jeżeli $\operatorname{red}(w)$
oznacza to słowo, działanie grupowe ma postać
\[
 u\cdot v=\operatorname{red}(uv),\qquad
 1=\varnothing,\qquad
 (x_1\cdots x_r)^{-1}=x_r^{-1}\cdots x_1^{-1}.
\]
\end{stwierdzenie}
\begin{proof}
Czytamy słowo od lewej do prawej i trzymamy
\emph{stos} już przeczytanych liter.
Gdy następna litera jest odwrotnością
litery na wierzchu stosu, zdejmujemy wierzch;
w przeciwnym razie dokładamy nową literę.
Po ostatnim kroku na stosie nie ma sąsiedniej
pary odwrotnej, więc otrzymaliśmy słowo
zredukowane.

Sprawdźmy niezależność od usuwania par.
Operacja dopisania litery $x$ do słowa
na stosie jest odwracalna: operacja
dopisania $x^{-1}$ cofa ją, zarówno gdy
pierwszy krok dołożył literę, jak i wtedy,
gdy ją skasował. Dlatego przeczytanie
fragmentu $xx^{-1}$ nie zmienia stanu stosu.
Usunięcie takiego fragmentu z dowolnego miejsca
słowa nie zmienia końcowego wyniku algorytmu.
Każde słowo równoważne daje zatem ten sam wynik.
Jeśli słowo jest już zredukowane, algorytm
pozostawia je bez zmian; w każdej klasie
może więc być tylko jedno takie słowo.

Sklejanie jest łączne przed redukcją, a
usuwanie par nie zmienia wyniku algorytmu.
Stąd $\operatorname{red}(\operatorname{red}(uv)w)
=\operatorname{red}(uvw)
=\operatorname{red}(u\operatorname{red}(vw))$.
Słowo puste jest jedynką. Zapisane w tezie
słowo odwrotne po sklejeniu z oryginałem
redukuje się kolejno od środka do słowa pustego.
To sprawdza wszystkie aksjomaty grupy.
\end{proof}

\begin{przyklad}[Słowa i działanie na słowach]
\label{prz:redukcja-slow-5}
W $F(\{a,b\})$ słowo
$ab\,b^{-1}a^{-1}ba$
redukujemy najpierw do
$aa^{-1}ba$, a następnie
do $ba$. Słowa $ab$ i $ba$
są już zredukowane i różne,
więc grupa ta nie jest przemienna.
Dla jednego generatora każde słowo
redukuje się do $a^m$ z jednoznacznym
$m\in\Z$; stąd $F(\{a\})\cong\Z$.

Grupa $F(A)$ działa na zbiorze zredukowanych
słów przez mnożenie z lewej:
$u\cdot w=\operatorname{red}(uw)$.
Na przykład $a\cdot(a^{-1}b)=b$.
Jest to działanie, bo łączność udowodniono
powyżej; odwrotna operacja do działania
przez $a$ to działanie przez $a^{-1}$.
Jest ono przechodnie, gdyż $w\cdot\varnothing=w$,
i wolne: jeśli $u\cdot w=w$, to w grupie
$uw=w$, zatem po pomnożeniu przez $w^{-1}$
dostajemy $u=1$. Ten sam zbiór ma również
działanie z prawej $w\cdot u=\operatorname{red}(wu)$.
Rysunek~\ref{fig:free-words-5} pokazuje
pojedyncze kroki tego drugiego działania.
Na przykład iloczyn dwóch zredukowanych słów
$ab$ i $b^{-1}a$ trzeba jeszcze skrócić:
$\operatorname{red}(ab\,b^{-1}a)=a^2$.
\end{przyklad}

\begin{figure}[htbp]
\centering
\begin{tikzpicture}[>=Latex,font=\small,
 every node/.style={circle,draw=manifoldblue!75!black,
 fill=manifoldblue!8,inner sep=2pt,minimum size=8mm}]
 \node (e) at (0,0) {$1$};
 \node (a) at (2,0) {$a$};
 \node (ai) at (-2,0) {$a^{-1}$};
 \node (b) at (0,1.5) {$b$};
 \node (bi) at (0,-1.5) {$b^{-1}$};
 \node (aa) at (3.8,.7) {$a^2$};
 \node (ab) at (3.8,-.7) {$ab$};
 \node (aiai) at (-3.8,.7) {$a^{-2}$};
 \node (aib) at (-3.8,-.7) {$a^{-1}b$};
 \node (bb) at (-.95,2.8) {$b^2$};
 \node (ba) at (.95,2.8) {$ba$};
 \node (bibi) at (-.95,-2.8) {$b^{-2}$};
 \node (bia) at (.95,-2.8) {$b^{-1}a$};
 \foreach \x/\y in {e/a,e/ai,e/b,e/bi,a/aa,a/ab,
 ai/aiai,ai/aib,b/bb,b/ba,bi/bibi,bi/bia}
   \draw[thick,manifoldblue!70!black] (\x)--(\y);
\end{tikzpicture}
\caption{Fragment grafu słów zredukowanych w
$F(\{a,b\})$. Krawędź oznacza dopisanie
z prawej jednej litery $a^{\pm1}$ albo $b^{\pm1}$
i ewentualną redukcję. Powrót po krawędzi
odpowiada dopisaniu litery odwrotnej.}
\label{fig:free-words-5}
\end{figure}

\begin{stwierdzenie}[Własność uniwersalna grupy wolnej]
\label{prop:free-universal-5}
Dla dowolnej grupy $K$ i dowolnego wyboru
elementów $h_a\in K$, $a\in A$, istnieje
dokładnie jeden homomorfizm $\Phi:F(A)\to K$
spełniający $\Phi(a)=h_a$.
\end{stwierdzenie}
\begin{proof}
Na słowach musi on mieć postać
\[
 \Phi(a_1^{\epsilon_1}\cdots a_r^{\epsilon_r})
 =h_{a_1}^{\epsilon_1}\cdots h_{a_r}^{\epsilon_r},
 \qquad \epsilon_i\in\{1,-1\}.
\]
Usunięcie sąsiedniej pary odwrotnej nie zmienia
iloczynu po prawej. Wzór zależy więc tylko
od klasy słowa i jest dobrze określony.
Sklejanie słów odpowiada mnożeniu obrazów,
zatem $\Phi$ jest homomorfizmem.
Każde słowo jest iloczynem generatorów
i ich odwrotności, co dowodzi jednoznaczności.
Właśnie brak dodatkowych relacji między
generatorami nazywamy \emph{wolnością}.
\end{proof}

W iloczynie wolnym $G*H$ działa podobny algorytm:
czytamy litery po kolei i trzymamy stos
liter różnych od jedynki, pochodzących naprzemiennie
z obu czynników. Nową literę z innego czynnika
dokładamy na stos; z tego samego mnożymy z literą
na wierzchu. Gdy iloczyn jest jedynką, usuwamy
wierzch i odsłaniamy poprzednią literę.
Operacja dopisania $x$ ma odwrotność: dopisanie
$x^{-1}$ z tego samego czynnika. Zatem usunięcie
sąsiednich liter jednego czynnika i zastąpienie
ich iloczynem nie zmienia końcowego stanu stosu.
Słowo już naprzemienne algorytm pozostawia bez zmian,
więc każda klasa ma dokładnie jedno takie słowo.
Stąd homomorfizmy $G\to K$ i $H\to K$
przedłużają się jednoznacznie na $G*H$.
To jest własność używana później w
Twierdzeniu~\ref{thm:van-kampen}.

Na przykład dla $G=\langle a\rangle\cong\Z$
i $H=\langle b\rangle\cong\Z$ własności
uniwersalne dają $G*H\cong F(\{a,b\})$.
Iloczyn słów naprzemiennych
$(ab)(b^{-1}a)$ redukuje się do $a^2$:
po skasowaniu $bb^{-1}$ zostaje $aa$.

\begin{definicja}[Prezentacja grupy]
\index{prezentacja grupy@Prezentacja grupy}
\label{def:prezentacja-grupy-5}
Jeśli $R$ jest zbiorem słów
w $F(A)$, to
\[
 \langle A\mid R\rangle
 :=F(A)/\langle\!\langle R\rangle\!\rangle
\]
nazywamy \emph{grupą o generatorach}
$A$ i \emph{relacjach} $R$.
Relacja $r=1$ oznacza, że jej
wszystkie konsekwencje po
mnożeniu i sprzęganiu również
stają się jedynką.
\end{definicja}

\begin{przyklad}[Prezentacja grupy torusa]
\label{prz:prezentacja-torusa-5}
Niech $F(a,b)$ będzie grupą wolną
na dwóch symbolach. Relacja
$aba^{-1}b^{-1}=1$
jest równoważna $ab=ba$.
Po jej narzuceniu każde słowo
można przestawić do postaci
$a^m b^n$, $m,n\in\Z$.
Odwzorowanie
$F(a,b)\to\Z^2$ dane przez
$a\mapsto(1,0)$,
$b\mapsto(0,1)$
zabija tę relację i wysyła
$a^m b^n$ na $(m,n)$.
Postacie te są więc wzajemnie
różne, a zatem
\[
 \langle a,b\mid aba^{-1}b^{-1}\rangle
 \cong\Z^2.
\]
Sklejenie przeciwnych boków kwadratu
tworzy torus. Brzeg kwadratu przed
sklejeniem odczytujemy jako słowo
$aba^{-1}b^{-1}$.
Topologiczne uzasadnienie, że
prezentacja opisuje $\pi_1(T^2)$,
podamy przy
Twierdzeniu~\ref{thm:van-kampen}.
\end{przyklad}

\section{Pierścienie, ideały i ilorazy}
\label{sec:pierscienie-idealy}
\index{pierscien@Pierścień}
\index{ideal@Ideał}

W grupie normalność podgrupy pozwala utworzyć
iloraz. Dla pierścieni podobną rolę pełni ideał:
trzeba zachować nie tylko dodawanie, lecz także
mnożenie klas przez dowolne elementy pierścienia.

\begin{definicja}[Pierścień przemienny i homomorfizm]
\label{def:pierscien-przemienny}
\emph{Pierścień przemienny z jedynką} $R$ ma
dwa działania $+$ i mnożenie. Grupa $(R,+,0)$
jest abelowa, mnożenie jest łączne i przemienne,
ma jedynkę $1$, a rozdzielność
$r(s+t)=rs+rt$ wiąże oba działania.
Pierścień zerowy z $1=0$ jest dozwolony.
\emph{Homomorfizm pierścieni}
$\varphi:R\to S$ zachowuje dodawanie, mnożenie
i jedynkę: $\varphi(1_R)=1_S$.
\end{definicja}

Przykładami są $\Z$, $\Q$ i pierścień
wielomianów $\mathbb K[t]$ nad ciałem $\mathbb K$.
W $\Z$ nie każdy niezerowy element ma odwrotność
mnożeniową: równanie $2x=1$ nie ma rozwiązania
w $\Z$. To odróżnia pierścień od ciała.

\begin{definicja}[Ideał i pierścień ilorazowy]
\label{def:ideal-iloraz}
\emph{Ideał} $I\subseteq R$ jest podgrupą
addytywną zamkniętą na mnożenie przez
elementy $R$: jeśli $r\in R$ i $x\in I$,
to $rx\in I$. Pierścień $R/I$ ma klasy
$[r]=r+I$ oraz działania
$[r]+[s]=[r+s]$, $[r][s]=[rs]$.
\end{definicja}

\begin{proof}[Dlaczego iloraz jest dobrze określony?]
Jeżeli $r'=r+i$ i $s'=s+j$ dla $i,j\in I$,
to $r'+s'-(r+s)=i+j\in I$ oraz
$r's'-rs=rj+is+ij\in I$.
Dlatego wynik obu działań nie zależy od
wyboru przedstawicieli. Jądro homomorfizmu
$\varphi:R\to S$ jest ideałem: dla
$\varphi(x)=0$ mamy
$\varphi(rx)=\varphi(r)\varphi(x)=0$.
Gdy $I=R$, iloraz jest pierścieniem zerowym
z $[1]=[0]$.
\end{proof}

\begin{przyklad}[Dwie relacje w pierścieniach]
\label{prz:idealy-ilorazy}
Niech $m\geq2$. Ideał $m\Z$ narzuca
relację $m=0$, a iloraz $\Z/m\Z$ jest pierścieniem
reszt modulo $m$. W $\mathbb K[t]$
ideał $(t):=t\mathbb K[t]=
\{tf(t):f\in\mathbb K[t]\}$ składa się z wielomianów
o zerowym wyrazie stałym. Ewaluacja
$f(t)\mapsto f(0)$ ma jądro $(t)$,
a odwzorowanie $[f]\mapsto f(0)$
daje izomorfizm
$\mathbb K[t]/(t)\cong\mathbb K$.
Istotnie, każda klasa zawiera dokładnie
jedną stałą, bo $f(t)-f(0)$ jest podzielny
przez $t$.
\end{przyklad}

\section{Moduły, homomorfizmy i ilorazy}
\label{sec:moduly}

Przez $R$ oznaczamy pierścień przemienny z jedynką
$1\ne0$ w sensie Definicji~\ref{def:pierscien-przemienny}.
Pozwala to traktować wszystkie moduły
jednakowo, bez rozróżniania strony działania.
Wiele definicji działa również dla pierścieni
nieprzemiennych i modułów lewych; zaznaczymy,
gdzie przemienność upraszcza zapis.

\begin{definicja}[Moduł nad pierścieniem]
\index{modul nad pierscieniem@Moduł nad pierścieniem}
\label{def:modul}
\emph{Moduł $M$ nad $R$} to grupa abelowa
$(M,+,0)$ z mnożeniem przez skalary
$R\times M\to M$, $(r,m)\mapsto rm$, takim że
dla $r,s\in R$ i $m,n\in M$ mamy
\[
 (r+s)m=rm+sm,\qquad r(m+n)=rm+rn,\qquad
 (rs)m=r(sm),\qquad 1m=m.
\]
Grupa abelowa oznacza tutaj, że dodawanie w $M$
jest łączne, przemienne, ma element $0$ i każdy
element $m$ ma element przeciwny $-m$.
\end{definicja}

Definicja różni się od definicji przestrzeni liniowej
jedynie tym, że nie wymagamy dzielenia przez
niezerowy skalar. Ta różnica ma poważne skutki:
moduł może nie mieć bazy, a równanie $rm=0$ może
zachodzić dla $r\ne0$ i $m\ne0$.
W przypadku grup abelowych, czyli $\Z$-modułów,
opisujemy to precyzyjnie następująco.

\begin{definicja}[Element torsyjny i $m$-torsja]
\label{def:torsja}
Element $x$ grupy abelowej $M$ jest
\emph{torsyjny}, jeśli istnieje liczba całkowita
$k\ne0$ taka, że $kx=0$. Zbiór takich elementów
oznaczamy $\operatorname{Tors}(M)$ i nazywamy
\emph{podgrupą torsyjną}. Dla ustalonego $m\geq2$
definiujemy \emph{$m$-torsję}
\[
 M[m]:=\{x\in M:mx=0\}=\ker(\times m:M\to M).
\]
Podgrupa $M[m]$ składa się z elementów zabijanych
przez tę jedną liczbę $m$; cały
$\operatorname{Tors}(M)$ dopuszcza różne niezerowe
liczby dla różnych elementów.
\end{definicja}

Oba zbiory są rzeczywiście podgrupami: jeśli
$kx=0$ i $\ell y=0$, to $k\ell(x-y)=0$;
dla ustalonego $m$ mamy od razu
$m(x-y)=mx-my=0$. Na przykład
$\operatorname{Tors}(\Z\oplus\Z/2)
=0\oplus\Z/2$, podczas gdy $\Z$
nie ma niezerowej torsji.

\begin{przyklad}[Od grup abelowych do spinorów]
\label{prz:moduly-podstawowe}
\begin{enumerate}[leftmargin=*]
\item Każda przestrzeń liniowa nad ciałem
      $\mathbb K$ jest $\mathbb K$-modułem.
\item Każda grupa abelowa $A$ jest $\Z$-modułem:
      $na$ dla $n>0$ oznacza sumę $n$ kopii $a$,
      $0a=0$, a $(-n)a=-(na)$. Aksjomaty modułu
      wynikają z praw dodawania. Odwrotnie,
      w każdym $\Z$-module aksjomaty wymuszają
      właśnie tę regułę, więc pojęcia są równoważne.
\item Dla $m\geq2$ grupa $\Z/m\Z$ jest $\Z$-modułem;
      niezerowa klasa $[1]$ spełnia $m[1]=0$. Nie ma bazy
      nad $\Z$. Istotnie, gdyby $[1]=\sum_i a_i e_i$ w takiej
      bazie, równość $\sum_i ma_i e_i=0$ i jednoznaczność
      współczynników wymuszałyby $ma_i=0$ w $\Z$ dla każdego $i$.
      Stąd wszystkie $a_i=0$, sprzecznie z $[1]\ne0$.
\item $R/I$ dla ideału $I\subseteq R$ jest
      $R$-modułem: $r[s]=[rs]$. Działanie jest
      dobrze określone, gdyż z $s-s'\in I$
      wynika $r(s-s')\in I$.
\item Moduł spinorów z Definicji~\ref{def:spinor}
      jest modułem nad zespoloną algebrą
      Clifforda. Tam algebra może być
      nieprzemienna i działanie jest \emph{lewe};
      obecny rozdział używa pierścienia
      przemiennego, a zatem uogólnia ideę
      działania algebry w kierunku
      przydatnym dla grup homologii.
\end{enumerate}
\end{przyklad}

\begin{definicja}[Podmoduł, homomorfizm, moduł wolny]
\index{podmodul homomorfizm modul wolny@Podmoduł, homomorfizm, moduł wolny}
\label{def:modul-map}
Podzbiór $N\subseteq M$ jest \emph{podmodułem},
jeśli zawiera $0$ i jest zamknięty na dodawanie,
branie elementów przeciwnych i mnożenie przez $R$.
\emph{Homomorfizm} $f:M\to N$ spełnia
$f(m+m')=f(m)+f(m')$ oraz $f(rm)=rf(m)$.
Piszemy $\operatorname{Hom}_R(M,N)$
na zbiór takich odwzorowań; przy dodawaniu punktowym
i działaniu $(rf)(m)=r f(m)$ jest on
$R$-modułem. Tutaj wykorzystujemy przemienność $R$:
$(rf)(sm)=rsf(m)=s(rf)(m)$, więc mnożenie homomorfizmu
przez skalar zachowuje jego $R$-liniowość.

Moduł \emph{wolny} z bazą $(e_i)_{i\in I}$
ma tę własność, że każdy jego element jest
\emph{jednoznaczną skończoną} sumą
$\sum_{i\in I}r_i e_i$. Modelowym przykładem
jest $R^{(I)}$, zbiór rodzin $(r_i)$
o skończonym nośniku.
Sumę prostą $M\oplus N$ tworzą pary $(m,n)$ z działaniami
składowymi; jej naturalne włączenia to $m\mapsto(m,0)$
i $n\mapsto(0,n)$, a rzuty wybierają odpowiednią współrzędną.
\end{definicja}

Na przykład $\operatorname{Hom}_{\Z}(\Z/2,\Z)=0$:
homomorfizm jest wyznaczony przez obraz $[1]$, lecz
$2f([1])=f(2[1])=0$ wymusza $f([1])=0$ w $\Z$.

To jest ten sam sens \emph{formalnych kombinacji},
którego użyliśmy w Rozdziale~\ref{ch:algebra-tensorowa}:
poza aksjomatami nie narzucamy żadnych relacji
między symbolami $e_i$. W szczególności
$\Z^k$ jest wolnym $\Z$-modułem, ale
$\Z/m\Z$ przy $m>1$ nie jest wolny.

Jeśli $N\subseteq M$ jest podmodułem, to
\emph{moduł ilorazowy} $M/N$ składa się z klas
$[m]=m+N$. Dwie klasy są równe dokładnie wtedy,
gdy $m-m'\in N$. Działania
$[m]+[m']=[m+m']$ i $r[m]=[rm]$
nie zależą od reprezentantów: zmianę
$m\mapsto m+n$ i $m'\mapsto m'+n'$
pochłania zamknięcie $N$ na $n+n'$ i $rn$.

\begin{twierdzenie}[Pierwsze twierdzenie o izomorfizmie]
\label{thm:modul-izomorfizm}
Dla homomorfizmu $f:M\to N$ jądro
$\ker f=\{m:f(m)=0\}$ i obraz
$\operatorname{im}f=\{f(m):m\in M\}$
są podmodułami, a wzór
\[
 \overline f:M/\ker f\longrightarrow\operatorname{im}f,
 \qquad [m]\longmapsto f(m)
\]
jest izomorfizmem $R$-modułów.
\end{twierdzenie}
\begin{proof}
Jeśli $f(m)=f(m')=0$, to
$f(m-m')=0$ i $f(rm)=r f(m)=0$;
stąd jądro jest podmodułem. Jeżeli
$x=f(m)$ i $y=f(m')$, to
$x+y=f(m+m')$ i $rx=f(rm)$, więc
obraz także jest podmodułem.
Jeżeli $[m]=[m']$, to $m-m'\in\ker f$
i $f(m)-f(m')=0$; wzór na $\overline f$
nie zależy więc od przedstawiciela.
Jest liniowy dzięki liniowości $f$.
Jest surjektywny na obraz z jego definicji.
Gdy $\overline f([m])=0$, mamy
$m\in\ker f$, a zatem $[m]=[0]$.
Jądro $\overline f$ jest zerowe, więc
odwzorowanie jest również iniektywne.
\end{proof}

\section{Iloczyn tensorowy modułów}
\label{sec:tensor-modulow-algebra}
\index{iloczyn tensorowy modulow@Iloczyn tensorowy modułów}

W Rozdziale~\ref{ch:algebra-tensorowa} tensorowaliśmy
przestrzenie liniowe nad ciałem. Nad pierścieniem
budujemy podobny obiekt, lecz nie wolno już korzystać
z bazy dowolnego modułu. Intuicja pozostaje ta sama:
symbol $m\otimes n$ ma być liniowy w obu argumentach,
a każda zbiliniowa reguła na parach powinna przechodzić
przez jedno odwzorowanie liniowe.

\begin{definicja}[Iloczyn tensorowy modułów]
\label{def:tensor-modulow}
Niech $M,N$ będą $R$-modułami. Weźmy wolny
$R$-moduł $F=R^{(M\times N)}$ na formalnych
symbolach $[m,n]$. Niech $K\subseteq F$ będzie
podmodułem generowanym przez wszystkie elementy
\[
\begin{aligned}
 &[m+m',n]-[m,n]-[m',n],\\
 &[m,n+n']-[m,n]-[m,n'],\\
 &[rm,n]-r[m,n],\qquad [m,rn]-r[m,n],
\end{aligned}
\]
gdzie $m,m'\in M$, $n,n'\in N$ i $r\in R$.
\emph{Iloczyn tensorowy} to $M\otimes_R N=F/K$,
a $m\otimes n$ oznacza klasę $[m,n]+K$.
\end{definicja}

Podmoduł generowany przez relacje składa się
z ich \emph{skończonych} kombinacji $R$-liniowych.
To jest algebraiczny sens narzucenia relacji.
Ukończony iloczyn tensorowy jest innym pojęciem:
wymaga dodatkowej topologii lub normy na modułach
i nie występuje w tej konstrukcji.
Z definicji
$rm\otimes n=m\otimes rn=r(m\otimes n)$.
Każdy element jest skończoną sumą tensorów
prostych, lecz nie musi być jednym takim tensorem.

\begin{twierdzenie}[Własność uniwersalna tensoru]
\label{thm:tensor-modulow-uniwersalnosc}
Dla każdego $R$-modułu $P$ i każdego
odwzorowania $R$-dwuliniowego
$b:M\times N\to P$ istnieje dokładnie
jeden homomorfizm $\widetilde b:M\otimes_R N\to P$
spełniający $\widetilde b(m\otimes n)=b(m,n)$.
\end{twierdzenie}
\begin{proof}
Na wolnym module $F$ istnieje dokładnie
jeden homomorfizm $L:F\to P$ z
$L([m,n])=b(m,n)$: określamy go na bazie
i przedłużamy liniowo. Addytywność i
$R$-liniowość $b$ w każdej zmiennej oznaczają,
że $L$ zeruje cztery rodzaje generatorów $K$.
Stąd $K\subseteq\ker L$, więc wzór
$\widetilde b(x+K)=L(x)$ jest dobrze określony.
Tensory proste rozpinają iloraz $F/K$;
wartości na nich wymuszają jednoznaczność.
Ten sam argument, zastosowany w obie strony,
pokazuje, że każda inna realizacja tej własności
jest kanonicznie izomorficzna z $M\otimes_R N$.
\end{proof}

\begin{przyklad}[Jednostka i iloraz przez ideał]
\label{prz:tensor-jednostka-ideal}
Odwzorowanie $R\otimes_R M\to M$,
$r\otimes m\mapsto rm$, jest izomorfizmem;
jego odwrotność wysyła $m$ na $1\otimes m$.
Jeśli $I\subseteq R$ jest ideałem, połóżmy
$IM=\{\sum_{j=1}^q i_jm_j:
i_j\in I,\ m_j\in M\}$.
Wówczas
\begin{equation}\label{eq:tensor-ideal-iloraz}
 (R/I)\otimes_R M\cong M/IM,
 \qquad (r+I)\otimes m\longmapsto rm+IM.
\end{equation}
Istotnie, wzór jest dwuliniowy i nie zależy
od wyboru $r$, bo $i m\in IM$ dla $i\in I$.
Odwrotność wysyła $m+IM$ na
$(1+I)\otimes m$. Jeśli $i\in I$, to
$(1+I)\otimes im=(i+I)\otimes m=0$;
znika więc na całym $IM$. Oba złożenia
są identycznościami na generatorach.
Dla $R=\Z$ i $I=m\Z$ otrzymujemy
$(\Z/m)\otimes_{\Z}A\cong A/mA$.
\end{przyklad}

\begin{stwierdzenie}[Tensorowanie odwzorowań]
\label{prop:tensor-maps-algebra}
Homomorfizm $f:M\to M'$ indukuje
$f\otimes\id_N:M\otimes_R N\to M'\otimes_R N$
przez $(f\otimes\id_N)(m\otimes n)=f(m)\otimes n$.
Jeśli $f$ jest surjekcją, to $f\otimes\id_N$
jest surjekcją. Tensorowanie nie musi
zachowywać iniekcji.
\end{stwierdzenie}
\begin{proof}
Zbiliniowość pary $(m,n)\mapsto f(m)\otimes n$
i Twierdzenie~\ref{thm:tensor-modulow-uniwersalnosc}
dają indukowany homomorfizm.
Gdy $f$ jest surjekcją, każdy tensor prosty
$m'\otimes n$ ma przeciwobraz $m\otimes n$
z $f(m)=m'$. Tensory proste rozpinają cel,
więc indukowany homomorfizm jest surjekcją.
Dla iniekcji $f=\times2:\Z\to\Z$ oraz
$N=\Z/2$ identyfikacja
$\Z\otimes_{\Z}\Z/2\cong\Z/2$ pokazuje,
że $f\otimes\id_N$ jest odwzorowaniem zerowym:
$(2z)\otimes[1]=z\otimes[2]=0$.
Nie jest zatem iniekcją.
\end{proof}

\begin{definicja}[Rozszerzenie skalarów]
\label{def:rozszerzenie-skalarow}
Jeśli $\varphi:R\to S$ jest homomorfizmem
pierścieni przemiennych z jedynką, to $S$
staje się $R$-modułem przez $r\cdot s=\varphi(r)s$.
Na $M\otimes_R S$ działa $S$ przez
$s\cdot(m\otimes a)=m\otimes sa$.
Ten $S$-moduł nazywamy \emph{rozszerzeniem
skalarów} modułu $M$ wzdłuż $\varphi$.
\end{definicja}

Sprawdźmy relację zbalansowania dla działania
$S$: dla $r\in R$ oraz $s,a\in S$ mamy
\[
 s\cdot((rm)\otimes a)
 =m\otimes\varphi(r)sa
 =m\otimes s\varphi(r)a
 =s\cdot(m\otimes\varphi(r)a).
\]
Przemienność $S$ uzasadnia środkową równość;
relacje addytywności sprawdza się składnikowo.
Dlatego działanie nie zależy od zapisu elementu
jako sumy tensorów prostych. Dla wolnego
$M=\bigoplus_{i\in J}Re_i$ mamy
\[
 M\otimes_R S\cong\bigoplus_{i\in J}Se_i,
 \quad
 (\textstyle\sum_i r_i e_i)\otimes s
 \longmapsto\textstyle\sum_i\varphi(r_i)s e_i.
\]
Odwrotność wysyła $\sum_i s_i e_i$ na
$\sum_i e_i\otimes s_i$; wszystkie sumy są
skończone. Podobnie tensorowanie zachowuje
sumy proste i izomorfizmy, lecz powyższy
przykład z $\times2$ ostrzega przed
wnioskowaniem o dowolnych iniekcjach.

\begin{uwaga}[Moduły projektywne]
\label{uwaga:moduly-projektywne}
Moduł $P$ jest \emph{projektywny}, gdy dla każdej
surjekcji $q:E\twoheadrightarrow B$ i każdego
homomorfizmu $g:P\to B$ istnieje podniesienie
$\widetilde g:P\to E$ z $q\widetilde g=g$.
Równoważnie, $P$ jest składnikiem prostym
pewnego modułu wolnego. Jeśli $P$ jest projektywny,
wybieramy wolny moduł $F$ na symbolach elementów
$P$ i jego naturalną surjekcję $q:F\twoheadrightarrow P$.
Podnosimy $\id_P$ przez $q$; otrzymane rozszczepienie
daje $F\cong P\oplus\ker q$. Odwrotnie,
odwzorowanie z wolnego modułu można podnieść,
wybierając przeciwobrazy wartości na bazie.
Jeśli $F=P\oplus Q$, to dla $g:P\to B$
podnosimy $g\circ\operatorname{pr}_P:F\to B$
i ograniczamy otrzymane odwzorowanie do $P$.
Projektywny nie musi być wolny: dla
$R=\Z/6$ ideał $3R=\{0,3\}$ jest składnikiem
prostym w $R=3R\oplus4R$, lecz jego dwa
elementy wykluczają izomorfizm z $R^{(J)}$.
\end{uwaga}

W zastosowaniach do nieprzemiennego pierścienia
grupowego $\Z[\pi]$ rozróżnia się moduł prawy
i lewy. Wtedy tensor tworzy się z prawego
modułu $M$ i lewego $N$, narzucając relację
$(m\lambda)\otimes n=m\otimes(\lambda n)$.
Wersja projektywności przez podnoszenie i
składniki proste wolnych działa również tam.
